\documentclass[11pt]{article}
\usepackage[T1]{fontenc}
\usepackage[a4paper,margin=25mm]{geometry}
\usepackage{amsmath,amssymb,amsthm,booktabs,array,longtable}
\usepackage[hidelinks]{hyperref}
\newtheorem{proposition}{Proposition}[section]
\newtheorem{corollary}[proposition]{Corollary}
\newcommand{\CT}{\operatorname{CT}}
\newcommand{\Dp}{D_p}
\title{Explicit Infinite Theta Expansion of the Somos--4 Kernel\\\large Nome-aligned formula with convolution-free and operator reductions}
\author{Thomas Scheuerle}
\date{23 September 2026}
\begin{document}
\maketitle

\begin{abstract}
Starting from the nome-aligned infinite theta expansion of the Somos--4 kernel, this revision removes the remaining convolution in the Fourier coefficients attached to $\sigma^2\wp$.  The antisymmetric coefficient of the kernel becomes an explicit product of the coefficients of $\theta_1^2$, a quadratic spectral factor, and one parity-dependent logarithmic derivative of theta constants.  Splitting the Fourier indices by parity then yields a compact differential-operator representation: the kernel is the sum of a wave-type Wronskian term and one two-component parity determinant.  In particular, the Gaussian constant in the sigma/theta normalization cancels from the antisymmetric coefficient, and the same-parity sectors vanish on both $m=\ell$ and $m=-\ell$.
\end{abstract}

\section{From the scaled generating function to the kernel}
Let
\[
G(x)=\sum_{k\ge0}g_kx^k,
\]
and let $H_n=H_n[G]$ denote the associated Hankel determinants. Under $G(x)\mapsto G(qx)$,
\[
H_n[G(qx)]=q^{n(n-1)}H_n[G].
\]
Define
\[
W_n:=H_nq^{n(n-1)},\qquad
S(t,q):=\sum_{n\in\mathbb Z}W_nt^n,
\]
and
\[
\boxed{\mathcal K(t,z;q_\ast):=S(tz,q_\ast)S(t/z,q_\ast).}
\]
Equivalently,
\[
\mathcal K(t,z;q_\ast)=\sum_{i,j\in\mathbb Z}W_iW_jt^{i+j}z^{i-j}.
\]
If
\[
H_{n+4}H_n=\alpha H_{n+3}H_{n+1}+\beta H_{n+2}^2,
\]
then
\[
[z^4]\mathcal K=\alpha q_\ast^6[z^2]\mathcal K+\beta q_\ast^8[z^0]\mathcal K.
\]

\section{Elliptic data and Fourier nome alignment}
Assume the standard nondegenerate sigma representation
\[
\boxed{H_n=A_{\rm S}B_{\rm S}^{\,n}
\frac{\sigma(z_0+n\kappa)}{\sigma(\kappa)^{n^2}}.}
\]
Choose fundamental half-periods $\omega_1,\omega_3$ and put
\[
\tau=\frac{\omega_3}{\omega_1},\qquad
p=e^{\pi i\tau},\qquad
\lambda=\frac{\pi}{2\omega_1},\qquad
a=\frac{\eta_1}{2\omega_1},
\]
where $\eta_1=\zeta(\omega_1)$. We use
\[
\boxed{\sigma(u)=C_\sigma e^{au^2}\theta_1(\lambda u,p)},
\qquad
\boxed{C_\sigma=\frac{2\omega_1}{\pi\theta_1'(0,p)}}.
\]
Define
\[
Q:=\frac{q}{\sigma(\kappa)},\qquad
\gamma:=\frac{\lambda\kappa}{2},\qquad
x_0:=\lambda z_0,
\]
and
\[
P:=Qe^{a\kappa^2}.
\]
For the infinite Fourier expansion we impose the alignment
\[
\boxed{P^2=p.}
\]
After choosing a branch of $p^{1/2}$, set
\[
\boxed{q_\ast:=\sigma(\kappa)p^{1/2}e^{-a\kappa^2}},
\qquad
\boxed{Q_\ast=p^{1/2}e^{-a\kappa^2},\quad P=p^{1/2}}.
\]
Finally put
\[
\boxed{X_\ast:=\frac{B_{\rm S}t}{q_\ast}e^{2az_0\kappa}.}
\]
This is specifically the Fourier alignment used below.  It is distinct, in general, from an alignment imposed after a finite torsion dissection.

\section{Parity theta functions}
For $\varepsilon\in\{0,1\}$ define
\[
\boxed{\Theta_\varepsilon(Y;p^{1/2})
:=\sum_{\substack{n\in\mathbb Z\\n\equiv\varepsilon\ ({\rm mod}\ 2)}}
p^{n^2/4}Y^n.}
\]
With
\[
\Phi_\rho(\xi):=\sum_{k\in\mathbb Z}\rho^{k^2}\xi^k
=f(\rho\xi,\rho/\xi),
\]
one has
\[
\boxed{\Theta_\varepsilon(Y;p^{1/2})=
p^{\varepsilon^2/4}Y^\varepsilon
\Phi_p(p^\varepsilon Y^2).}
\]

\section{Fourier coefficients and elimination of the convolution}
Define $a_\ell$ and $b_\ell$ by
\[
\theta_1(x,p)^2=\sum_{\ell\in\mathbb Z}a_\ell(p)e^{2i\ell x},
\]
\[
\lambda^2\bigl(\theta_1'(x,p)^2-\theta_1(x,p)\theta_1''(x,p)\bigr)
-2a\theta_1(x,p)^2
=\sum_{\ell\in\mathbb Z}b_\ell(p)e^{2i\ell x}.
\]
The first sequence has the parity formulas
\[
\boxed{a_{2j}=p^{2j^2}\theta_2(0,p^2)},
\qquad
\boxed{a_{2j+1}=-p^{(2j+1)^2/2}\theta_3(0,p^2).}
\]
Equivalently,
\[
a_\ell=(-1)^\ell\sum_{n\in\mathbb Z}
p^{(n+1/2)^2+(\ell-n-1/2)^2}.
\]
Let
\[
\boxed{\Dp:=p\frac{d}{dp}}.
\]

\begin{proposition}[Convolution-free formula for $b_\ell$]
For every $\ell\in\mathbb Z$,
\[
\boxed{b_\ell=4\lambda^2(\ell^2-\Dp)a_\ell-2a\,a_\ell.}
\]
\end{proposition}

\begin{proof}
The direct convolution is
\[
b_\ell=2\lambda^2(-1)^\ell\sum_{n\in\mathbb Z}
(2n+1)(2\ell-2n-1)
p^{(n+1/2)^2+(\ell-n-1/2)^2}-2a\,a_\ell.
\]
Set $r=2n+1-\ell$. Then
\[
(2n+1)(2\ell-2n-1)=\ell^2-r^2,
\]
\[
(n+\tfrac12)^2+(\ell-n-\tfrac12)^2=\frac{\ell^2+r^2}{2},
\]
and $r\equiv1-\ell\pmod2$. Hence
\[
a_\ell=(-1)^\ell
\sum_{\substack{r\in\mathbb Z\\r\equiv1-\ell\ ({\rm mod}\ 2)}}
p^{(\ell^2+r^2)/2}.
\]
Applying $\Dp$ and eliminating the $r^2$-weighted sum gives
\[
(-1)^\ell\sum_r(\ell^2-r^2)p^{(\ell^2+r^2)/2}
=2(\ell^2-\Dp)a_\ell,
\]
which proves the claim.
\end{proof}

Let $\pi(\ell)\in\{0,1\}$ denote the numerical parity of $\ell$, and put
\[
T_0(p):=\theta_2(0,p^2),\qquad T_1(p):=\theta_3(0,p^2).
\]
Then
\[
\boxed{b_\ell=a_\ell\left[
2\lambda^2\ell^2-4\lambda^2\Dp\log T_{\pi(\ell)}(p)-2a
\right].}
\]
Thus $b_\ell$ contains no convolution and no data beyond the elliptic lattice.

\section{Antisymmetric coefficient and its cancellations}
Define the single parity modulus
\[
\boxed{\Delta(p):=\Dp\log\frac{\theta_2(0,p^2)}{\theta_3(0,p^2)}.}
\]

\begin{proposition}[Explicit antisymmetric coefficient]
For all $\ell,m\in\mathbb Z$,
\[
\boxed{
a_\ell b_m-b_\ell a_m
=2\lambda^2a_\ell a_m
\left[m^2-\ell^2+2\Delta(p)\bigl(\pi(m)-\pi(\ell)\bigr)\right].}
\]
\end{proposition}

\begin{proof}
Insert the preceding formula for $b_k$. The terms $-2a\,a_k$ cancel identically. The difference of the two parity logarithmic derivatives is $-\Delta$, $0$, or $\Delta$ according to the parities of $(\ell,m)$, which gives the displayed formula.
\end{proof}

The four parity sectors are therefore
\[
\begin{array}{c|c}
(\pi(\ell),\pi(m))&
(a_\ell b_m-b_\ell a_m)/(2\lambda^2a_\ell a_m)\\ \hline
(0,0)&m^2-\ell^2\\
(1,1)&m^2-\ell^2\\
(0,1)&m^2-\ell^2+2\Delta(p)\\
(1,0)&m^2-\ell^2-2\Delta(p)
\end{array}
\]
In particular, every same-parity term contains
\[
m^2-\ell^2=(m-\ell)(m+\ell),
\]
so the same-parity contribution vanishes on both the diagonal $m=\ell$ and the anti-diagonal $m=-\ell$. The full coefficient always vanishes on $m=\ell$ by antisymmetry.

\section{Convolution-free infinite-theta kernel}
The rank-two identity is
\[
\mathcal K(t,z;q_\ast)=A_{\rm S}^2\sum_{\varepsilon=0}^1
\left(\mathcal R_\varepsilon\widehat{\mathcal D}_\varepsilon
-\widehat{\mathcal R}_\varepsilon\mathcal D_\varepsilon\right).
\]
Substitution of the explicit antisymmetric coefficient gives
\[
\boxed{
\begin{aligned}
\mathcal K(t,z;q_\ast)
={}&2\lambda^2A_{\rm S}^2C_\sigma^4e^{2az_0^2}
\sum_{\varepsilon=0}^1\sum_{\ell,m\in\mathbb Z}
e^{2i\ell x_0}a_\ell a_m\\
&\times\left[m^2-\ell^2
+2\Delta(p)\bigl(\pi(m)-\pi(\ell)\bigr)\right]\\
&\times\Theta_\varepsilon(X_\ast e^{2i\ell\gamma};p^{1/2})
\Theta_\varepsilon(ze^{2im\gamma};p^{1/2}).
\end{aligned}}
\]
Expanding the two parity theta functions gives the entirely Ramanujan form
\[
\boxed{
\begin{aligned}
\mathcal K(t,z;q_\ast)
={}&2\lambda^2A_{\rm S}^2C_\sigma^4e^{2az_0^2}
\sum_{\varepsilon=0}^1\sum_{\ell,m\in\mathbb Z}
e^{2i\ell x_0}a_\ell a_m\\
&\times\left[m^2-\ell^2
+2\Delta(p)\bigl(\pi(m)-\pi(\ell)\bigr)\right]\\
&\times p^{\varepsilon^2/2}
(X_\ast e^{2i\ell\gamma})^\varepsilon
(ze^{2im\gamma})^\varepsilon\\
&\times\Phi_p\!\left(p^\varepsilon X_\ast^2e^{4i\ell\gamma}\right)
\Phi_p\!\left(p^\varepsilon z^2e^{4im\gamma}\right).
\end{aligned}}
\]

\section{A two-component differential-operator reduction}
The parity split can be reorganized further. For $r,s\in\{0,1\}$ define the theta lifts
\[
\mathcal R_{\varepsilon,r}(x;t):=
\sum_{\substack{\ell\in\mathbb Z\\\ell\equiv r\ ({\rm mod}\ 2)}}
a_\ell e^{2i\ell x}
\Theta_\varepsilon(X_\ast e^{2i\ell\gamma};p^{1/2}),
\]
\[
\mathcal D_{\varepsilon,s}(y;z):=
\sum_{\substack{m\in\mathbb Z\\m\equiv s\ ({\rm mod}\ 2)}}
a_m e^{2imy}
\Theta_\varepsilon(ze^{2im\gamma};p^{1/2}).
\]
For functions of $x$ and $y$, set
\[
\boxed{\mathcal W(F,G):=\frac14\left((\partial_x^2F)G-F(\partial_y^2G)\right).}
\]
On a Fourier monomial this operator has symbol $m^2-\ell^2$. Put
\[
\mathcal R_{\varepsilon,+}:=\mathcal R_{\varepsilon,0}+\mathcal R_{\varepsilon,1},
\qquad
\mathcal D_{\varepsilon,+}:=\mathcal D_{\varepsilon,0}+\mathcal D_{\varepsilon,1}.
\]

\begin{proposition}[Wave term plus parity determinant]
The full kernel is
\[
\boxed{
\begin{aligned}
\mathcal K(t,z;q_\ast)
={}&2\lambda^2A_{\rm S}^2C_\sigma^4e^{2az_0^2}
\sum_{\varepsilon=0}^1
\Bigl\{
\mathcal W(\mathcal R_{\varepsilon,+},\mathcal D_{\varepsilon,+})\\
&\qquad\qquad
+2\Delta(p)
\bigl(
\mathcal R_{\varepsilon,0}\mathcal D_{\varepsilon,1}
-\mathcal R_{\varepsilon,1}\mathcal D_{\varepsilon,0}
\bigr)
\Bigr\}_{x=x_0,\,y=0}.
\end{aligned}}
\]
\end{proposition}

\begin{proof}
The operator $\mathcal W$ supplies $m^2-\ell^2$ in all four parity sectors. The determinant
\[
\det\begin{pmatrix}
\mathcal R_{\varepsilon,0}&\mathcal R_{\varepsilon,1}\\
\mathcal D_{\varepsilon,0}&\mathcal D_{\varepsilon,1}
\end{pmatrix}
=
\mathcal R_{\varepsilon,0}\mathcal D_{\varepsilon,1}
-\mathcal R_{\varepsilon,1}\mathcal D_{\varepsilon,0}
\]
adds $+2\Delta$ in the even--odd sector, $-2\Delta$ in the odd--even sector, and zero in the same-parity sectors.
\end{proof}

This representation gives a structural interpretation of the double series. Its main part is a wave-type differential image of two separable theta lifts. All residual coupling between the Fourier parities is carried by one $2\times2$ determinant and the single lattice scalar $\Delta(p)$.

\section{Free and derived quantities}
\renewcommand{\arraystretch}{1.2}
\begin{longtable}{>{\raggedright\arraybackslash}p{0.18\textwidth}>{\raggedright\arraybackslash}p{0.22\textwidth}>{\raggedright\arraybackslash}p{0.50\textwidth}}
\toprule
Symbol & Status & Meaning / determination \\
\midrule
$G(x)$ & input & Original moment generating function. Its Hankel determinants give $H_n$.\\
$q_\ast$ & Fourier alignment & $q_\ast=\sigma(\kappa)p^{1/2}e^{-a\kappa^2}$ after choosing a branch of $p^{1/2}$.\\
$A_{\rm S},B_{\rm S},z_0,\kappa$ & orbit data & Gauge constants, initial elliptic phase, and elliptic translation step.\\
$\omega_1,\omega_3$ & lattice data & Fundamental half-periods; $\tau=\omega_3/\omega_1$.\\
$p,\lambda,a,C_\sigma$ & derived lattice data & Defined by the sigma/theta normalization above.\\
$P$ & aligned value & $P=p^{1/2}$, so the outer bilateral theta functions have nome $P^2=p$.\\
$\gamma,x_0,X_\ast$ & derived & Shift phase, initial phase, and the combined $t$-variable.\\
$a_\ell$ & explicit coefficient & Fourier coefficient of $\theta_1^2$, with closed parity formulas.\\
$b_\ell$ & eliminated coefficient & $b_\ell=4\lambda^2(\ell^2-\Dp)a_\ell-2aa_\ell$; no convolution is needed.\\
$\Delta(p)$ & derived lattice scalar & Logarithmic nome derivative of $\theta_2(0,p^2)/\theta_3(0,p^2)$; it is the only non-polynomial coefficient in the parity coupling.\\
\bottomrule
\end{longtable}

\section{Analytic qualification and verification status}
The rearrangements may be interpreted formally whenever all sums are handled coefficientwise in a valid formal setting. For an analytic interpretation, assume $|p|<1$ and impose a domain in the variables for which the bilateral sums and all differentiated sums converge absolutely and locally uniformly. Under such hypotheses, application of $\Dp$, $\partial_x^2$, and $\partial_y^2$ term by term is justified.

The convolution elimination and the operator formula are algebraic consequences of the displayed Fourier series. They can be checked coefficientwise by truncating the theta series and comparing:
\begin{enumerate}
\item the original convolution for $b_\ell$ with $4\lambda^2(\ell^2-\Dp)a_\ell-2aa_\ell$;
\item the original determinant $a_\ell b_m-b_\ell a_m$ with the parity formula;
\item each of the four parity sectors with the wave-plus-determinant operator.
\end{enumerate}

\section{Conclusion}
The infinite theta expansion now has a convolution-free coefficient structure. The sequence $b_\ell$ is generated from $a_\ell$ by a second-order spectral operator and a nome derivative. Antisymmetrization removes the Gaussian constant and leaves only the quadratic factor $m^2-\ell^2$ together with one parity modulus $\Delta(p)$. Finally, the complete double series is a wave-type bilinear operator plus a two-component parity determinant. This is a substantive simplification of the kernel and provides a more promising starting point for modular identities, special-value reductions, and numerical verification.

\end{document}
